%============================================================
% HALAMAN SAMPUL DEPAN (COVER)
%   - Judul: uppercase bold, piramida terbalik (atur dengan \\)
%   - Teks putih "HALAMAN SAMPUL" di antara judul dan "Skripsi"
%   - Logo: tinggi 4 cm, lebar 5,5 cm, terpusat
%   - Identitas: nama 12pt bold, NIM, prodi, fakultas, kampus, tahun
% Halaman ini dihitung halaman i (romawi) tetapi nomornya tidak
% ditampilkan.
%============================================================
\thispagestyle{empty}
\onehalfspacing

\begin{center}
  % ── Judul skripsi (piramida terbalik) ──
  \vspace*{1em}
  {\bfseries\MakeUppercase{\JudulSkripsi}\par}
  \JudulStudiKasus

  % ── Teks tersembunyi putih: "HALAMAN SAMPUL" ──
  \HalamanAwal{HALAMAN SAMPUL}

  \vspace{0.4em}
  {\bfseries Skripsi\par}

  % ── Logo kampus: tinggi 4 cm, lebar maksimal 5,5 cm ──
  \vfill
  \includegraphics[height=4cm,width=5.5cm,keepaspectratio]{assets/amikom/logo_amikom.png}

  % ── Identitas mahasiswa ──
  \vfill
  Disusun oleh

  \vspace{0.3em}
  {\bfseries \NamaMahasiswa\par}
  {\bfseries \NIM\par}

  % ── Program studi, fakultas, universitas, kota, tahun ujian ──
  \vfill
  {\bfseries PROGRAM STUDI \MakeUppercase{\NamaProdi}\par}
  {\bfseries \MakeUppercase{\NamaFakultas}\par}
  {\bfseries UNIVERSITAS AMIKOM PURWOKERTO\par}
  {\bfseries PURWOKERTO\par}
  \vspace{0.4em}
  {\bfseries \TahunUjian\par}
\end{center}

\clearpage
